% Generated draft: requires mathematical and source review.
\documentclass[11pt]{article}
\usepackage[margin=25mm]{geometry}
\usepackage{fontspec}
\setmainfont{Latin Modern Roman}
\usepackage{amsmath,amssymb,mathtools,unicode-math}
\setmathfont{Latin Modern Math}
\usepackage{enumitem,needspace,xurl,hyperref}
\setlength{\parindent}{0pt}
\setlength{\parskip}{5pt}
\setlength{\emergencystretch}{3em}
\newcommand{\N}{\mathbb{N}}
\newcommand{\Z}{\mathbb{Z}}
\newcommand{\Q}{\mathbb{Q}}
\newcommand{\R}{\mathbb{R}}
\newcommand{\C}{\mathbb{C}}
\newcommand{\F}{\mathbb{F}}
\newcommand{\A}{\mathbb{A}}
\newcommand{\PP}{\mathbb{P}}
\newcommand{\E}{\mathbb{E}}
\newcommand{\eps}{\varepsilon}
\newcommand{\dd}{\,\mathrm{d}}
\begin{document}
\section*{UnsolvedMath ID: 20000329}
\textbf{Source problem code:} catalogue entry
\subsection{Definitions and mathematical statement}
Let $U$ be smooth and \'etale over a moduli stack $\mathcal M_H$ of polarized manifolds, and let $Y$ be a log compactification with boundary $\Gamma=Y\setminus U$. Under the usual log-smooth convention, $Y$ is smooth projective and $\Gamma$ a reduced simple-normal-crossings divisor. The logarithmic cotangent bundle consists of forms with logarithmic poles, with determinant $\omega_Y(\Gamma)$.

Viehweg's original questions are whether (1) $\Omega_Y^1(\log\Gamma)$ is weakly positive with respect to $U$, and (2) $\omega_Y(\Gamma)$ is ample with respect to $U$. One usual weak-positivity convention requires that, for every ample line bundle $A$ and integer $a>0$, some $b>0$ makes $\operatorname{Sym}^{ab}E\otimes A^b$ generated by global sections over $U$. The source does not define its convention or ``ample with respect to $U$''; that phrase must not be silently replaced by ampleness merely of the restriction to $U$ or by global ampleness on $Y$. In common linear-system versions, a positive power must define a suitable embedding over $U$, but the precise requirement needs the source setup. The ground field, moduli data $H$, and class of polarized manifolds are also unspecified.

The source reports both statements known for period domains and says (2) implies (1), without hypotheses or references. It suggests configuration-space analogues as potentially more tractable; these are a related program, not automatically \'etale charts of the stated moduli stack.

The supplied research report treats ordered two-point configurations on a smooth projective curve $C$: set $Y=C\times C$, $\Gamma=\Delta_C$, and $U=Y\setminus\Delta_C$. For positive genus it claims $\Omega_Y^1(\log\Delta_C)$ nef, hence Viehweg-weakly positive over all of $Y$, and $\omega_Y(\Delta_C)$ ample exactly when $g(C)\geq2$. Nefness means nonnegative degree of the tautological quotient line along every curve in the bundle's projectivization.

For an elliptic curve with an origin, it reports that each graph $\{(x,x+a):x\in C\}$ for $a\ne0$ is a complete curve in $U$ contracted by every complete-linear-system map of every positive power of the log canonical bundle. It uses this to obstruct ampleness with respect to $U$, not just global ampleness. The report explicitly leaves the original moduli-stack problem open. Its configuration claims and positivity implications remain unreviewed.
\subsection{Short English statement}
Ask for logarithmic cotangent weak positivity and log canonical ampleness over an étale moduli chart. The report studies the suggested two-point configuration analogue, giving positive-genus nefness and an elliptic interior-curve obstruction. That analogue does not settle the original moduli question or define its relative positivity conventions.
\subsection{Sources}
\begin{itemize}
\item Local UnsolvedMath registry, permanent ID $20000329$, source code \texttt{catalogue entry}; created and updated 2026-08-14. Status \texttt{partially\_solved}; classification \texttt{PARTIAL-PROGRESS}; difficulty L3; proposer field null, with source attribution to Viehweg.
\item AIM Problem Lists, Algebraic geometry, workshop \emph{Compact moduli spaces and birational geometry}, Problem 2.2, source type PDF, \url{https://aimath.org/WWN/birational/birational.pdf}. Canonical location \texttt{aim-algebraic-geometry-notes.json}, \texttt{notes[328]}, tag \texttt{problem}; statement status \texttt{corrected\_verified}. Verification reports official-PDF correction of duplicated numbering, a stray cotangent-notation symbol, spacing, and $\mathcal M_H$. Original remarks and literature field empty.
\item The original supplies period-domain and implication claims but no exact theorem references. The report adds the curve configuration results without exact primary bibliography or ground-field conventions.
\item Supplied research attempt 1: result \texttt{partial\_result}, run status \texttt{open}, rigor label \texttt{proved}, low novelty confidence. Candidate contribution \texttt{configuration\_space\_proposition\_and\_obstruction}: positive-genus log-cotangent nefness, genus-at-least-two log canonical ampleness, and elliptic translation-graph obstruction. Explicit caveat: machine-generated research aid, not peer reviewed; independently verify before citation.
\item Dataset attribution: Ulam AI / UnsolvedMath Contributors, CC BY 4.0; original problem attribution remains to the supplied AIM workshop collection.
\end{itemize}
\end{document}
